\documentclass[11pt,a4paper]{article}

\usepackage[T1]{fontenc}
\usepackage[utf8]{inputenc}
\usepackage[margin=22mm]{geometry}
\usepackage{hyperref}

\setlength{\parindent}{0pt}
\setlength{\parskip}{0.6em}
\newcommand{\code}[1]{\texttt{\detokenize{#1}}}

\title{Running FeatFloWer Heat and Gendie Cases\\from an Installed Tree}
\author{FeatFloWer workflow handout}
\date{\today}

\begin{document}
\maketitle

\section{The installation, input, and runtime case}

The launchers distinguish three locations:

\begin{description}
  \item[Installation directory] Contains executables, Python launchers,
  partitioner code, and shipped defaults. It is read-only during a normal run.
  Typical locations are \code{INSTALL/bin/q2p1_gendie} and
  \code{INSTALL/bin/heat}.

  \item[Input directory] Contains the configuration and geometry that define a
  simulation. It is selected with \code{-f}.

  \item[Runtime case directory] Contains staged input, generated meshes,
  partition data, logs, dumps, and visualization output. It is selected with
  \code{-C}. The solver uses this directory as its working directory.
\end{description}

Because launcher and input paths may be absolute, both applications can be
started from any current working directory.

\section{Standalone heat application}

\subsection{Heat input contract}

The directory passed with \code{-f} must initially contain:

\begin{verbatim}
heat_input/
|-- heat.s3d                 required heat configuration
|-- sampleRigidBody.xml      required boundary description
|-- steel.off                OFF files referenced by heat.s3d
|-- titan.off
|-- melt.off
|-- wire_0.off
|-- ...
|-- wall_1.off               OFF files referenced by sampleRigidBody.xml
|-- ...
|-- wall_30.off
`-- meshDir/                 optional fallback if meshing cannot generate one
    |-- file.prj
    |-- Mesh.tri
    |-- wall_1.par
    `-- ...
\end{verbatim}

The input \code{sampleRigidBody.xml} is placed beside \code{heat.s3d}, not in
an input \code{start} directory. On every launch it is copied to
\code{<case>/start/sampleRigidBody.xml}. This input copy is authoritative and
replaces an older runtime copy.

All top-level \code{*.off} and \code{*.OFF} files are copied from the input
directory to the runtime case root. Thus entries such as
\code{steel.off} in \code{heat.s3d} and \code{wall_1.off} in
\code{sampleRigidBody.xml} resolve correctly. Geometry stored below an input
subdirectory is not copied and must instead use a path that is valid relative
to the runtime case, or an absolute path.

The optional input \code{meshDir} is a complete fallback mesh. Every
parametrization file named by its \code{file.prj} must be present. During a run
the active mesh is located below \code{<case>/_data/meshDir}.

\subsection{Heat execution}

Example with separate input and runtime directories:

\begin{verbatim}
python3 /opt/FeatFloWer/INSTALL/bin/heat/heat_start.py \
  -C /scratch/user/heat_01/RUN \
  -f /scratch/user/heat_01/CASE \
  -n 64
\end{verbatim}

The launcher creates the runtime directory skeleton, stages
\code{heat.s3d}, \code{sampleRigidBody.xml}, and top-level OFF files, runs
\code{s3d_mesher}, partitions for \code{n-1} worker ranks, and launches
\code{heat} with \code{n} total MPI ranks. Rank zero is the control rank, so
at least two ranks are required. Add \code{--use-srun} on a suitable Slurm
allocation.

Important generated paths include:

\begin{verbatim}
RUN/_data/heat.s3d
RUN/_data/meshDir/
RUN/start/sampleRigidBody.xml
RUN/mesh_names.offs
RUN/_vtk/
RUN/_dump/
\end{verbatim}

\subsection{Installed heat test case}

The installation contains a complete heat example at:

\begin{verbatim}
<install>/bin/heat/_ianus/HEAT
\end{verbatim}

It includes \code{heat.s3d}, \code{sampleRigidBody.xml}, all heat and wall OFF
geometry, and a complete fallback \code{meshDir}. It can be run directly as
the input while keeping output in a separate runtime case:

\begin{verbatim}
python3 /opt/FeatFloWer/INSTALL/bin/heat/heat_start.py \
  -C /scratch/user/heat_example/RUN \
  -f /opt/FeatFloWer/INSTALL/bin/heat/_ianus/HEAT \
  -n 64
\end{verbatim}

\section{Gendie through q2p1\_sse}

\subsection{Gendie input contract}

A typical YAML-driven gendie runtime case contains:

\begin{verbatim}
gendie_case/
|-- e3d_input/
|   |-- setup.e3d or Extrud3D.dat
|   |-- *.off                  referenced geometry
|   `-- meshDir/               optional input mesh
|-- plan.yaml                  simulation plan
`-- _data_BU/                  optional case-specific parameter overrides
    `-- q2p1_paramV_DIE_test.dat
\end{verbatim}

The installation provides the solver executables, mesher, partitioner,
standard YAML plans, and default \code{_data_BU} templates. Missing runtime
directories and static runtime defaults are seeded automatically. A file in
the case directory overrides the corresponding installed template.

The YAML plan defines ordered \code{init}, \code{main}, and \code{final}
stages. Supported solver steps include momentum, temperature, and material
distribution, depending on the installed applications.

\subsection{Gendie execution}

\begin{verbatim}
python3 /opt/FeatFloWer/INSTALL/bin/q2p1_gendie/e3d_start_yaml.py \
  -C /scratch/user/gendie_01 \
  -f /scratch/user/gendie_01/e3d_input \
  -y /scratch/user/gendie_01/plan.yaml \
  -n 64 -x
\end{verbatim}

Here \code{-x} requests the gendie short-test behavior. It is not a heat
launcher option. As with heat, \code{-C} is the runtime working directory and
the command may be issued from anywhere.

\section{Module-aware build and case setup helper}

The repository helper
\code{cmake_cheat_setup_from_Raphael.sh} configures a CGAL-enabled build,
builds and installs the application tree, prepares an example input, and
prints (but does not execute) the simulation command.

The module selection is mandatory. For gendie, use:

\begin{verbatim}
./cmake_cheat_setup_from_Raphael.sh \
  -M gendie /scratch/user/gendie_01
\end{verbatim}

Invoking the helper without \code{-M gendie} or \code{-M heat} stops before
configuration, compilation, or installation.

By default, the installed RH\_LOC example is copied to
\code{/scratch/user/gendie_01/e3d_input}. An alternative template can be
selected with \code{-i INPUT}.

For heat, the installed \code{_ianus/HEAT} test case is used by default:

\begin{verbatim}
./cmake_cheat_setup_from_Raphael.sh \
  -M heat /scratch/user/heat_01/RUN
\end{verbatim}

It is copied to \code{RUN/heat_input}. To use a custom heat input template,
add \code{-i}:

\begin{verbatim}
./cmake_cheat_setup_from_Raphael.sh \
  -M heat \
  -i /scratch/user/heat_01/CASE \
  /scratch/user/heat_01/RUN
\end{verbatim}

The printed run command uses \code{-C RUN} together with
\code{-f RUN/heat_input}. If that input directory is already populated, it is
preserved when \code{-i} is omitted.

Useful environment overrides are:

\begin{verbatim}
INSTALL_PREFIX=/path/to/INSTALL
BUILD_DIR=/path/to/build
BUILD_JOBS=8
MPI_RANKS=64
INPUT_SOURCE=/path/to/input
\end{verbatim}

For example:

\begin{verbatim}
INSTALL_PREFIX=$PWD/../INSTALL-HEAT BUILD_JOBS=8 MPI_RANKS=64 \
  ./cmake_cheat_setup_from_Raphael.sh \
  -M heat -i /scratch/user/heat_01/CASE /scratch/user/heat_01/RUN
\end{verbatim}

\section{Quick consistency checks}

For heat, the number of XML boundary shapes should match the generated wall
surfaces and boundary IDs:

\begin{verbatim}
grep -c '<BoundaryShape ' CASE/sampleRigidBody.xml
ls CASE/wall_*.off | wc -l
\end{verbatim}

After staging, verify that every XML mesh reference resolves from the runtime
case directory:

\begin{verbatim}
cd RUN
grep -o 'meshFile="[^"]*"' start/sampleRigidBody.xml |
sed 's/meshFile="//; s/"//' |
while IFS= read -r file; do
  test -f "$file" || echo "MISSING: $file"
done
\end{verbatim}

If heat prints all \code{wall_*.par} names and then crashes in
\code{projectonboundaryid}, first check that the runtime XML contains the
matching \code{BoundaryDescription} and that all referenced \code{wall_*.off}
files exist in the runtime case.

\end{document}
